\section{Relations and closures}
\label{sec:strings-05}

\begin{definition}[Reflexive transitive closure]
Let \(R\subseteq X\times X\).  The \emph{reflexive transitive closure} of
\(R\), written \(R^*\), is the least reflexive and transitive relation on
\(X\) that contains \(R\).
\end{definition}

\begin{lemma}[Path characterization of the reflexive transitive closure]
\label{lem:rtc-path-characterization}
Let \(R\subseteq X\times X\) and let \(x,y\in X\).  Then \(xR^*y\) iff there
are \(n\in\N\) and \(x_0,\ldots,x_n\in X\) such that \(x_0=x\), \(x_n=y\),
and \(x_iRx_{i+1}\) for every \(i<n\).
\textit{Status: complete proof.}
\end{lemma}
\begin{proof}
Call a finite sequence \(x_0,\ldots,x_n\) with the displayed three properties
a path from \(x\) to \(y\) along \(R\), and let \(S\subseteq X\times X\) be
the relation holding between \(x\) and \(y\) iff such a path exists.  The
claim is \(R^*=S\), and each inclusion is proved separately.

For \(R^*\subseteq S\), it suffices to show that \(S\) is reflexive and
transitive and contains \(R\), since \(R^*\) is the least such relation.
Reflexivity: for any \(x\), take \(n=0\) and \(x_0=x\); the condition
\(x_iRx_{i+1}\) for every \(i<0\) holds vacuously.  Containment of \(R\): if
\(xRy\), take \(n=1\), \(x_0=x\), \(x_1=y\).  Transitivity: let
\(x_0,\ldots,x_m\) be a path from \(x\) to \(y\) and \(y_0,\ldots,y_n\) a
path from \(y\) to \(z\).  Define \(u_0,\ldots,u_{m+n}\) by \(u_i=x_i\) for
\(i\leq m\) and \(u_{m+j}=y_j\) for \(j\leq n\); the two clauses agree at
the one shared index \(m\), since \(x_m=y=y_0\).  Then \(u_0=x\),
\(u_{m+n}=z\), and for every \(i<m+n\): if \(i<m\), then
\(u_iRu_{i+1}\) by the first path; if \(m\leq i<m+n\), write \(i=m+j\) with
\(j<n\), and then \(u_i=y_jRy_{j+1}=u_{i+1}\) by the second path.  So
\(u_0,\ldots,u_{m+n}\) is a path from \(x\) to \(z\).

For \(S\subseteq R^*\), induction on \(n\) proves that every path
\(x_0,\ldots,x_n\) from \(x\) to \(y\) along \(R\) gives \(xR^*y\).  Base
case \(n=0\): then \(y=x_0=x\), and \(xR^*x\) because \(R^*\) is reflexive.
Inductive step: assume the claim for \(n\), and let \(x_0,\ldots,x_{n+1}\)
be a path from \(x\) to \(y\).  The initial segment \(x_0,\ldots,x_n\) is a
path from \(x\) to \(x_n\), so \(xR^*x_n\) by the induction hypothesis.
Also \(x_nRx_{n+1}\) and \(R\subseteq R^*\), so \(x_nR^*y\).  Transitivity
of \(R^*\) gives \(xR^*y\).
\end{proof}

\begin{definition}[Equivalence closure]
The \emph{equivalence closure} of a relation \(R\subseteq X\times X\) is the
least equivalence relation on \(X\) containing \(R\).
\end{definition}

These closures will be used for reduction sequences, proof normalization, and
quotient constructions.  They must not be confused: a reflexive transitive
closure need not be symmetric.

